%=====================================================================
\chapter{Managing Orders, Customers and Refunds}\label{ch:orders}
%=====================================================================
\locator{5}{Part V --- Running and Extending the Shop}

\begin{outcomes}
By the end of this chapter you will be able to:
\begin{tight}
  \item \textbf{Explain} how the admin panel changes an order, and \textbf{say}
        why it does not do so directly.
  \item \textbf{Predict} which status changes move stock, from the rule rather
        than from the status names.
  \item \textbf{Distinguish} the two things the platform calls a refund, and
        \textbf{state} what each one actually does.
  \item \textbf{Identify} the work a refund requires that no screen performs.
  \item \textbf{Design} a back-office procedure that keeps money, stock and
        records in agreement.
\end{tight}
\end{outcomes}

\begin{prereq}
\chref{checkout}: order statuses, the history table, and status 0.
\chref{dashboard}: that the dashboard's figures are filtered differently.
\chref{gateway}: reconciliation, and that the shop's record and the provider's
must be made to agree.
\end{prereq}

\begin{keyterms}
back office $\bullet$ order history $\bullet$ status transition $\bullet$ stock
subtraction $\bullet$ return $\bullet$ return action $\bullet$ refund $\bullet$
store credit $\bullet$ customer transaction $\bullet$ API signature $\bullet$
allow-list $\bullet$ audit trail $\bullet$ reconciliation
\end{keyterms}

\section{The Admin Panel Does Not Edit Orders}

The first surprise of the back office is architectural. When an administrator
changes an order's status, the admin panel does not write to the database. It
\textbf{signs a request and calls the storefront's own API over HTTP}, and the
storefront makes the change.

\begin{conceptbox}[Why that is the right design]
Three things must happen when an order's status changes: stock must move,
the customer may need an email, and a history row must be written. Those rules
are not simple, and they must be identical whether the change comes from a
human clicking a dropdown, from a payment provider's callback
(\chref{gateway}), or from a script.

\smallskip
Writing them once, in the storefront, and having everything else call in is
the only arrangement in which they cannot drift apart. The alternative ---
the admin panel having its own copy --- is how a shop ends up with orders
whose stock was never adjusted because they were edited from the wrong screen.

\smallskip
\textbf{The consequence nobody mentions: the admin panel cannot change an order
unless the API accepts the call.} The back office is a client of the shop, and
it can be locked out of it.
\end{conceptbox}

\begin{measured}[Out of the box, the API refuses every call]
Produced by running \texttt{m18\_orders.py --api} in \texttt{code/}. The
request is signed exactly as the startup controller verifies it: an HMAC-SHA1
over
the route, the call, the username, the host, the store, the language, the
currency, an MD5 of the request body and a timestamp.

\smallskip
\begin{tabular}{@{}L{5.6cm}L{8.2cm}@{}}
Rows in \texttt{oc\_api\_ip} & \textbf{0} \\
A correctly signed call & \texttt{HTTP 403 --- You do not have permission} \\
After adding one address & accepted \\
\end{tabular}

\smallskip
\textbf{The API user ships enabled, and its address allow-list ships empty.} An
empty allow-list permits nothing --- the check is ``is the caller's address in
this list'', and the answer for an empty list is always no. So every signed
call is refused, including the admin panel's own.

\smallskip
This is the third fault of this shape in the course, after \chref{checkout}'s
missing setting and \chref{security}'s unread \texttt{.htaccess}: a feature
that is present, configured, enabled, and inert. In each case the system
answered, politely, and did nothing.
\end{measured}

\section{Stock Moves on Transitions, Not on Statuses}

The rule that governs stock is not written on any screen, and it is not what
the status names suggest.

\begin{definitionbox}[The rule, from \texttt{catalog/model/checkout/order.php}]
The statuses are partitioned into a set --- the union of
\texttt{config\_processing\_status} and \texttt{config\_complete\_status} ---
and everything else.

\smallskip
\textbf{Moving \emph{into} that set subtracts stock. Moving \emph{out of} it
adds stock back. Moving within it does nothing.}

\smallskip
In this shop the set is $\{1, 2, 3, 5, 12\}$ --- Pending, Processing, Shipped,
Complete and \emph{Reversed}. Everything else, including 0, Canceled, Denied,
Failed and \emph{Refunded}, is outside it.
\end{definitionbox}

\begin{measured}[One order, seven status changes, and the stock after each]
From \texttt{code/m18\_orders.py --stock}: one order for one unit of product~43
and two of product~40, driven through the API.

\rowcolors{2}{white}{lightbg}
\begin{longtable}{@{}L{3.4cm}L{6.0cm}L{2.0cm}L{2.0cm}@{}}
\toprule
\rowcolor{primary!15}
\textbf{Status} & \textbf{Where it sits} & \textbf{Stock 43} &
\textbf{Stock 40} \\
\midrule
\endhead
(before the order) & --- & 904 & 920 \\
1 Pending & in the set & 903 & 918 \\
3 Shipped & still in the set & 903 & 918 \\
7 Canceled & \textbf{outside} & \textbf{904} & \textbf{920} \\
1 Pending & back inside & 903 & 918 \\
\textbf{11 Refunded} & \textbf{outside} & \textbf{904} & \textbf{920} \\
1 Pending & back inside & 903 & 918 \\
\textbf{12 Reversed} & \textbf{in the set} & \textbf{903} & \textbf{918} \\
\bottomrule
\end{longtable}

\smallskip
\textbf{Shipped changed nothing}, although it is the status at which the goods
physically leave. The stock left when the order was paid for.

\smallskip
\textbf{Refunded and Reversed behave oppositely.} Both names mean money went
back to the customer. Status~11 returns the goods to stock; status~12 does not,
because somebody put 12 in \texttt{config\_processing\_status}. Nothing on any
screen says so, and the two statuses sit next to each other in the dropdown.

\smallskip
\textbf{This is a setting, so it is your problem.} Whoever configures a shop
decides which statuses are ``processing'', and in doing so decides the stock
behaviour of every status change the shop will ever make. It is one of the
highest-consequence settings in the admin panel and it looks like a list of
checkboxes.
\end{measured}

\section{Two Different Things Are Called a Refund}

\begin{alertbox}[A return marked ``Refunded'' does nothing at all]
The platform has a returns system: \adminpath{Sales, Returns}, with its own
statuses (Pending, Awaiting Products, Complete), five reasons, and four
actions --- \uiel{Waiting for return}, \uiel{Credit Issued},
\uiel{Replacement Sent} and \uiel{Refunded}.

\smallskip
Reading \texttt{admin/model/sale/returns.php}, recording a return performs
exactly two statements:

\begin{tight}
  \item \texttt{UPDATE oc\_return SET return\_status\_id = \dots}
  \item \texttt{INSERT INTO oc\_return\_history \dots}
\end{tight}

\smallskip
\textbf{And nothing else.} It does not change the order's status. It does not
restore stock. It does not create a customer transaction or any store credit.
It does not move a single rupee.

\smallskip
\textbf{A return marked ``Refunded'' is a note.} It records that somebody
decided to refund. The money must be sent back by hand through the payment
provider, the stock adjusted by hand, and the order's own status changed by
hand --- and if none of that is done, the order still reads \uiel{Complete} to
everyone who looks at it, including the dashboard and every report.
\end{alertbox}

\dsfig{%
\begin{tikzpicture}[
  act/.style={rounded corners=3pt, draw=#1, line width=1.0pt, fill=#1!8,
              text=#1!60!black, font=\bfseries\scriptsize, align=center,
              text width=2.9cm, minimum height=0.9cm},
  done/.style={font=\scriptsize\bfseries, text=greenhl},
  todo/.style={font=\scriptsize\bfseries, text=accent},
  note/.style={font=\scriptsize, text=neutral, align=left, text width=5.8cm,
               anchor=west}
]
  \node[act=secondary] (r) at (0,3.0) {Mark the return\\Refunded};
  \node[act=accent]    (m) at (0,1.9) {Send the money back};
  \node[act=accent]    (s) at (0,0.8) {Put the goods back};
  \node[act=accent]    (o) at (0,-0.3) {Change the order's\\status};
  \node[act=accent]    (c) at (0,-1.5) {Reverse the\\affiliate commission};

  \node[note, text=greenhl] at (1.9,3.0)
       {\faIcon{check}~The platform does this. Two SQL statements.};
  \node[note, text=accent] at (1.9,1.9)
       {\faIcon{times}~By hand, in the provider's own console.};
  \node[note, text=accent] at (1.9,0.8)
       {\faIcon{times}~By hand, or by moving the order to status 11 --- but
        not 12.};
  \node[note, text=accent] at (1.9,-0.3)
       {\faIcon{times}~By hand, or every report still counts the sale.};
  \node[note, text=accent] at (1.9,-1.5)
       {\faIcon{times}~Never, unless you built it (\chref{discounts}).};
\end{tikzpicture}}%
{What a refund actually requires, and the one part of it the software performs.
Four of the five are somebody's responsibility, and in most small shops nobody
has been told which.}{refundwork}

\begin{conceptbox}[The procedure a shop actually needs]
Because the software performs one fifth of a refund, the other four fifths have
to be a written procedure, in this order:

\begin{tightnum}
  \item \textbf{Send the money back first}, through the provider, and record
        the provider's reference.
  \item \textbf{Move the order to status 11}, not 12, which returns the goods
        to stock --- and only if the goods actually came back.
  \item \textbf{Record the return} with the provider's reference in the
        comment, so that the shop's record and the provider's can be matched
        later (\chref{gateway}).
  \item \textbf{Reverse any commission}, which nothing will do for you.
  \item \textbf{Check the dashboard agrees} the next morning. The gap between
        what the shop thinks it sold and what the provider says it collected is
        the only thing that catches a step somebody forgot.
\end{tightnum}

\smallskip
\textbf{Steps 2 and 3 are in that order on purpose.} A return recorded before
the order is moved leaves a window in which the shop's records say ``refunded''
and its stock and totals say ``sold''. The window is usually a few minutes, and
month-end is when somebody finds the one that lasted a year.
\end{conceptbox}

\section{Customers, and What You Are Keeping}

The customer side of the back office is shorter but carries more risk, because
the records are about people.

\begin{examplebox}[What \adminpath{Customers, Customers} lets you do, and the
question attached to each]
\begin{tight}
  \item \textbf{Edit a customer's details.} Who is allowed to, and is the change
        recorded? An order copies its addresses (\chref{checkout}), so editing
        the customer does not rewrite history --- which is correct, and worth
        confirming rather than assuming.
  \item \textbf{Reset a password.} The shop cannot show you the old one, and
        anyone who offers to has stored it wrongly (\chref{accounts}).
  \item \textbf{See everything they bought.} Convenient for support, and a
        complete purchase history of a named person --- which
        \chref{legal} treats as the regulated asset it is.
  \item \textbf{Add store credit} (\texttt{oc\_customer\_transaction}). This is
        real money the shop owes, and it is the one place in the back office
        where an administrator can create a liability with no approval step.
  \item \textbf{Delete a customer.} Their orders remain, because the orders are
        the shop's accounting records and are not the customer's to delete. Be
        able to explain that distinction before you are asked it.
\end{tight}
\end{examplebox}

\begin{pitfall}
\begin{tight}
  \item \textbf{Assuming the admin panel writes to the database.} It calls the
        storefront, and can be refused.
  \item \textbf{Leaving \texttt{oc\_api\_ip} empty} and then reporting that
        order editing is broken without knowing why.
  \item \textbf{Reading the status names to predict stock.} Refunded returns
        stock and Reversed does not.
  \item \textbf{Putting a money-returned status into
        \texttt{config\_processing\_status}}, which is how Reversed got there.
  \item \textbf{Believing a return marked Refunded refunded anything.}
  \item \textbf{Recording the return before moving the order}, leaving a window
        where the two disagree.
  \item \textbf{Refunding without reversing commission}, which
        \chref{discounts} named as the commonest unreconciled figure in a small
        shop.
  \item \textbf{Editing an order to ``correct'' it} instead of changing its
        status, which loses the history a dispute would be settled from.
\end{tight}
\end{pitfall}

\begin{lab}[Refund an order properly, then find what the shop still believes]
\begin{lstlisting}[language=bash]
# --- 1. Find out why order editing does not work --------------------
cd code
docker compose exec db mariadb -ushop -pshoppw opencart \
  -e "SELECT * FROM oc_api; SELECT * FROM oc_api_ip;"
# One table has a row and the other does not. Say what that means.

# --- 2. Prove it from outside ---------------------------------------
python m18_orders.py --api

# --- 3. Predict the stock table, then watch it ----------------------
# Which of these returns goods to stock: Shipped, Canceled, Refunded,
# Reversed? Commit to an answer for each BEFORE running:
python m18_orders.py --stock

# --- 4. Explain the two that surprised you --------------------------
docker compose exec db mariadb -ushop -pshoppw opencart \
  -e "SELECT value FROM oc_setting \
      WHERE \`key\` IN ('config_processing_status', \
                       'config_complete_status');"

# --- 5. Record a return in the admin panel --------------------------
# Admin: Sales > Returns > Add New. Action: Refunded. Save.
# Now check: did the order's status change? Did stock move? Did any
# money move? Check all three in the database, not on the screen.
python m18_orders.py --refund

# --- 6. Do the refund properly --------------------------------------
# Follow the five-step procedure in this chapter. Write down which
# steps the software performed and which you performed.

# --- 7. Find the disagreement ---------------------------------------
# Admin: Reports > Sales > Orders. Does the total match what the shop
# actually kept? Account for every difference.

# --- 8. Write the procedure -----------------------------------------
# One page a shop assistant could follow, with the order of steps
# and a stated reason for that order.

# --- 9. Clean up ----------------------------------------------------
python m18_orders.py --clean
\end{lstlisting}

\smallskip
\textbf{What to notice.} Step 5 is the one students argue about. The screen says
\uiel{Refunded}, the save succeeds, a confirmation appears --- and the database
shows two rows written in a table nothing else reads.
\end{lab}

\begin{checkpoint}
\begin{tightnum}
  \item An order is moved from Pending to Shipped, then to Reversed. State what
        happened to stock at each step and why, and say what you would change in
        the shop's configuration if you wanted Reversed to behave like Refunded.
  \item A shop records 40 returns a month, all marked \uiel{Refunded}, and its
        stock figures drift steadily away from reality. Explain the mechanism
        precisely, say what else is also drifting, and give the smallest change
        that would stop it.
  \item The admin panel reports that it cannot change an order. Describe how you
        would diagnose this in four steps, naming what you would look at each
        time, and explain why the fault is easy to misread as a permissions
        problem with the administrator's account.
\end{tightnum}
\end{checkpoint}

\begin{chaptersummary}
\begin{tight}
  \item The admin panel changes an order by signing a request and calling the
        storefront's API, so that stock, notification and history follow one
        code path whatever triggered the change.
  \item Measured: the API user ships enabled with an \emph{empty} address
        allow-list, which permits nothing --- so every signed call is refused,
        including the admin panel's own. The third inert-feature fault in this
        course.
  \item Stock is decided by the transition between sets of statuses, not by the
        status. Moving into the processing-or-complete set subtracts; moving
        out of it adds back; moving within it does nothing.
  \item Measured: Shipped moved no stock. Canceled and Refunded returned it.
        \textbf{Reversed did not}, because 12 sits in
        \texttt{config\_processing\_status} --- a setting, and therefore your
        responsibility.
  \item A return marked \uiel{Refunded} performs two SQL statements and nothing
        else: no order status change, no stock, no credit, no money. It is a
        note.
  \item A refund is five jobs and the software does one. The other four need a
        written procedure, in a stated order, because the order of steps decides
        how long the records disagree.
  \item Store credit is a liability an administrator can create with no approval
        step. Deleting a customer does not delete their orders, and should not.
\end{tight}
\end{chaptersummary}

\begin{reviewq}
  \item \textbf{(MCQ)} Moving an order to status 12 (Reversed) causes stock to:
        (a)~be returned; (b)~be subtracted again; (c)~stay as it is;
        (d)~depend on the payment method.
  \item \textbf{(MCQ)} Recording a return with the action \uiel{Refunded}:
        (a)~refunds the customer; (b)~restores stock; (c)~changes the order's
        status; (d)~writes two rows and does nothing else.
  \item \textbf{(Short)} Explain why the admin panel calls the storefront's API
        instead of writing to the database, and give one problem that design
        prevents.
  \item \textbf{(Short)} Why does an empty \texttt{oc\_api\_ip} refuse every
        call rather than allowing every call? Give the general lesson about
        allow-lists.
  \item \textbf{(Short)} An order's addresses are copied at checkout. Explain
        what this means for editing a customer's address afterwards, and why it
        is the correct behaviour.
  \item \textbf{(Applied)} A shop's stock counts have drifted by several hundred
        units over six months. Give four distinct mechanisms from this chapter
        that could produce that, say how you would test for each in the
        database, and state which you would expect to find first.
  \item \textbf{(Applied)} Write the refund procedure for a shop taking 200
        orders a day with a 6\% return rate, two staff in the back office and a
        payment provider whose console is separate from the shop. State each
        step, who performs it, what is recorded, the order and the reason for
        the order, and the daily check that would catch a missed step.
\end{reviewq}
